% ------------------------------------------------------------------ packages
\usepackage[T1]{fontenc}
\usepackage{lmodern}
\usepackage[margin=1.1in]{geometry}
\usepackage{amsmath,amssymb,amsthm,mathtools}
\usepackage{microtype}
\usepackage{enumitem}
\usepackage{array,booktabs,tabularx}
\usepackage{xcolor}
\usepackage{tikz}
\usetikzlibrary{arrows.meta,positioning,calc,fit,backgrounds,shapes.geometric,patterns,decorations.pathreplacing}
\usepackage[most]{tcolorbox}
\usepackage[font=small,labelfont=bf,skip=4pt]{caption}
\usepackage[colorlinks=true,linkcolor=blue!50!black,citecolor=teal!55!black,
            urlcolor=violet!65!black,bookmarksnumbered=true]{hyperref}
\usepackage{cleveref}

% ------------------------------------------------------------------ colours
\definecolor{ideacol}{RGB}{255,247,224}
\definecolor{ideaframe}{RGB}{214,178,80}
\definecolor{warncol}{RGB}{253,235,235}
\definecolor{warnframe}{RGB}{200,90,90}
\definecolor{checkcol}{RGB}{235,243,252}
\definecolor{checkframe}{RGB}{90,130,190}
\definecolor{plancol}{RGB}{240,247,240}
\definecolor{planframe}{RGB}{90,150,100}
\definecolor{contractcol}{RGB}{246,242,250}
\definecolor{contractframe}{RGB}{130,100,170}
\definecolor{graytext}{RGB}{90,90,90}

% ------------------------------------------------------------------ theorems
% thmtools is used so that cleveref prints the right environment name even
% though all statement types share one counter (and therefore one numbering).
\usepackage{thmtools}
\declaretheorem[name=Theorem,numberwithin=section,style=definition,
  refname={Theorem,Theorems},Refname={Theorem,Theorems}]{theorem}
\declaretheorem[name=Proposition,numberlike=theorem,style=definition,
  refname={Proposition,Propositions},Refname={Proposition,Propositions}]{proposition}
\declaretheorem[name=Lemma,numberlike=theorem,style=definition,
  refname={Lemma,Lemmas},Refname={Lemma,Lemmas}]{lemma}
\declaretheorem[name=Corollary,numberlike=theorem,style=definition,
  refname={Corollary,Corollaries},Refname={Corollary,Corollaries}]{corollary}
\declaretheorem[name=Definition,numberlike=theorem,style=definition,
  refname={Definition,Definitions},Refname={Definition,Definitions}]{definition}
\declaretheorem[name=Example,numberlike=theorem,style=definition,
  refname={Example,Examples},Refname={Example,Examples}]{example}
\declaretheorem[name=Construction,numberlike=theorem,style=definition,
  refname={Construction,Constructions},Refname={Construction,Constructions}]{construction}
\declaretheorem[name=Remark,numberlike=theorem,style=remark,
  refname={Remark,Remarks},Refname={Remark,Remarks}]{remark}

\crefname{figure}{Figure}{Figures}
\crefname{table}{Table}{Tables}
\crefname{section}{Section}{Sections}
\crefname{appendix}{Appendix}{Appendices}

\numberwithin{equation}{section}

% ------------------------------------------------------------------ boxes
% "In words" gloss, attached above a formal statement
\newtcolorbox{words}{
  enhanced, breakable=false, colback=black!4, colframe=black!4, boxrule=0pt,
  leftrule=2.4pt, arc=0pt, outer arc=0pt, left=7pt, right=7pt, top=3pt, bottom=3pt,
  fontupper=\small\itshape}

% section contract
\newtcolorbox{contract}[1][]{
  enhanced, breakable=true, colback=contractcol, colframe=contractframe,
  boxrule=0.6pt, arc=2pt, left=7pt, right=7pt, top=13pt, bottom=5pt,
  fonttitle=\small\bfseries\sffamily, title={What this section does},
  attach boxed title to top left={xshift=8pt, yshift=-8pt},
  boxed title style={colback=contractframe, colframe=contractframe, arc=1pt},
  #1}

% proof plan
\newtcolorbox{planbox}[1][Proof plan]{
  enhanced, breakable=true, colback=plancol, colframe=planframe,
  boxrule=0.5pt, arc=2pt, left=7pt, right=7pt, top=12pt, bottom=4pt,
  fonttitle=\small\bfseries, title={#1},
  attach boxed title to top left={xshift=6pt, yshift=-7pt},
  boxed title style={colback=planframe, colframe=planframe, arc=1pt, left=3pt,right=3pt,top=1pt,bottom=1pt}}

% idea / warning / checkpoint
\newtcolorbox{ideabox}[1][]{
  enhanced, breakable=true, colback=ideacol, colframe=ideaframe, boxrule=0.6pt,
  arc=2pt, left=7pt, right=7pt, top=12pt, bottom=4pt,
  fonttitle=\small\bfseries, title={Idea},
  attach boxed title to top left={xshift=6pt, yshift=-7pt},
  boxed title style={colback=ideaframe, colframe=ideaframe, arc=1pt, left=3pt,right=3pt,top=1pt,bottom=1pt},
  #1}
\newtcolorbox{warnbox}[1][]{
  enhanced, breakable=true, colback=warncol, colframe=warnframe, boxrule=0.6pt,
  arc=2pt, left=7pt, right=7pt, top=12pt, bottom=4pt,
  fonttitle=\small\bfseries, title={Watch out},
  attach boxed title to top left={xshift=6pt, yshift=-7pt},
  boxed title style={colback=warnframe, colframe=warnframe, arc=1pt, left=3pt,right=3pt,top=1pt,bottom=1pt},
  #1}
\newtcolorbox{checkpoint}[1][]{
  enhanced, breakable=true, colback=checkcol, colframe=checkframe, boxrule=0.6pt,
  arc=2pt, left=7pt, right=7pt, top=12pt, bottom=4pt,
  fonttitle=\small\bfseries, title={Where we are},
  attach boxed title to top left={xshift=6pt, yshift=-7pt},
  boxed title style={colback=checkframe, colframe=checkframe, arc=1pt, left=3pt,right=3pt,top=1pt,bottom=1pt},
  #1}

% "used in" line
\newcommand{\usedin}[1]{\par\smallskip\noindent{\small\color{graytext}\textbf{Used in:} #1.}}

% ------------------------------------------------------------------ macros
\newcommand{\abs}[1]{\lvert #1\rvert}
\newcommand{\card}[1]{\lvert #1\rvert}
\DeclareMathOperator{\rank}{rank}
\DeclareMathOperator{\im}{im}
\DeclareMathOperator{\Span}{span}
\newcommand{\R}{\mathbb R}
\newcommand{\from}{\leftarrow}
\newcommand{\disjoint}{\sqcup}
\newcommand{\Ne}{N_1^+}
\newcommand{\Nee}{N_2^+}
\newcommand{\Nminus}{N_1^-}
\newcommand{\restr}{\upharpoonright}
\DeclareMathOperator{\dom}{dom}
\DeclareMathOperator{\coker}{coker}
\newcommand{\Emph}[1]{\emph{#1}}
\newcommand{\code}[1]{\texttt{#1}}

% theorem-attached "In words" convenience
\newcommand{\inwords}[1]{\begin{words}#1\end{words}}

% step heading inside proofs
\newcommand{\steplabel}[1]{\par\smallskip\noindent\textbf{#1}\quad}
